\section{Data and reference trees}
\label{sec:data}

We use WikiArt-Refined \citep{tan2019wikiart}, $\sim$81{,}446
paintings labelled with one of 27 styles, with the supplied 70/30
train/val split. The corpus is heavily class-imbalanced
($133.3\times$ ratio between Impressionism with 13{,}060 examples and
Analytical Cubism with 77; full distribution in
Appendix~\ref{app:data}). Two design choices follow: the prototype
classifier has no per-class bias (only geometric distance enters the
logit), and we report balanced accuracy alongside top-1.

\paragraph{Reference trees.} WikiArt ships no hierarchy and
art-history offers several credible taxonomies, so we use three. The
\emph{default} tree (\cref{app:tree}) follows standard
art-historical lineage. The \emph{chronological} tree groups styles
into six era buckets (Renaissance, Baroque--Rococo, 19th century,
early 20th century, modern post-war, non-Western). The \emph{flat}
tree makes every style a direct child of the root and serves as a
deliberate null whose pairwise distances are constant off-diagonal.
For a tree $T$ with leaves $u, v$, tree distance is the unweighted
shortest-path edge count
$d_T(u, v) = \mathrm{depth}(u) + \mathrm{depth}(v) - 2\,\mathrm{depth}(\mathrm{LCA}(u, v))$.
We additionally build two \emph{empirical} reference trees by
agglomerative clustering of class-mean encoder features
(\cref{ssec:local}): one from CLIP and one from DINOv2, used only
for evaluation.
